\documentclass[UTF8,11pt,a4paper]{ctexart}
\usepackage[margin=2.4cm]{geometry}
\usepackage{booktabs,longtable,enumitem,hyperref,xcolor}
\hypersetup{colorlinks=true,urlcolor=blue,linkcolor=black}
\setlist{nosep}
\title{AI 工具使用详情}
\author{}
\date{2026年8月30日}
\begin{document}
\pagestyle{plain}
\maketitle

\section{工具名称、版本或型号}
本候选成果使用 OpenAI Codex 桌面版（GPT-5 系列编程研究智能体）。平台在本地任务工件中未暴露精确服务端构建号，因此不虚构具体小版本。运行环境为 Windows、Python 3.12，代码求解依赖见支撑材料中的 \texttt{requirements.txt}。

\section{具体用途和环节}
AI 工具用于：读取与解释官方题面和数据字典；提出可检验的数学模型；编写、调试并运行问题一时序集成预测、问题二混合整数规划、问题三泊松联合仿真和动态资源规划、问题四官方赛程采集与比较代码；生成图表、论文候选稿与机械审计脚本。AI 未获得隐藏测试集的真实转播观看人数，也未执行任何现实赛事操作。

\section{主要提示方式与使用过程}
核心提示要求在独立冻结目录内完成四问，不得读取历史 C 题成果；官方输入只读；所有新工件写入本项目；输出严格匹配模板；生成可复现代码、论文、AI 声明和支撑包。初稿在核心审查门前停止；独立审查返回后实施一次实质回修，当前停在交付审查前，不得自行标记 READY。

典型交互方式不是要求 AI 直接给出结论，而是要求其持续执行“题面核对--数据审计--可运行求解--不变量检查--交付渲染”。遇到求解耗时，明确要求采用有界时间，保存当前最好可行解，并披露最优性缺口。

\section{采纳、修改和核验}
\begin{longtable}{p{2.7cm}p{5.3cm}p{5.3cm}}
\toprule
环节 & AI 产出 & 已实施的程序化核验\\
\midrule
问题一 & 赛前特征与岭回归、极端随机树、随机森林、梯度提升候选模型 & 固定日期顺序的五折阻断验证；显式排除赛后变量；检查140/72行与单位；披露集成权重和误差同用OOF、并非嵌套外层验证\\
问题二 & 场馆--时段 MILP & 风险按官方$[0,1]$原值，完整成本含17.95百万美元恒定启用成本；显式加入96条60小时约束并重算全部硬约束；3.0570\% gap只属于轮次日期预筛模型\\
问题三 & 20,000 次联合泊松仿真与静/动态资源 MILP & 晋级权重和为32；票务、转播、成本共享更新后Min--Max边界，吸引力和风险用原值；逐日重算资源与预算；两个MILP gap=0\\
问题四 & FIFA API 采集与跨赛事比较 & 保存原始响应与URL；保留11项诊断；另核对Q2八项、共同边界和官方权重总分，并以81组代理敏感性检出排序翻转\\
交付 & CSV、论文 PDF、图表和支撑包 & 模板表头逐字匹配、ID集合、有限值、PDF渲染、页数、匿名性和 SHA-256 检查\\
\bottomrule
\end{longtable}

\section{人工复核状态与责任边界}
本文件随候选交付生成。参赛队仍须在正式提交前逐项人工复核题意解释、模型选择、代码、数值、引用和论文表述，并把队号“000”替换为真实队号。当前成果不得因 AI 或程序运行成功而自动视为科学结论正确；核心审查回修不等于交付审查通过，也不得自行标记为 READY。

\end{document}
